\subsection{Data Source and Sample Period}

The analysis uses five-minute candlestick data for the
BTCUSDT USD-margined perpetual futures contract obtained from
the Binance Public Data repository
\citep{binancepublicdata}. The sample covers January 1, 2022
through December 31, 2025, with all timestamps recorded in
Coordinated Universal Time (UTC).

The 48 monthly files contain 420,768 observations. Each record
includes opening, high, low, and closing prices, trading
volume, quote volume, the number of trades, and taker-buy
volume, following the Binance Kline data structure
\citep{binancepublicdata}. The five-minute frequency provides
sufficient intraday detail for constructing hourly volatility
and activity measures while remaining computationally
manageable.

Observations are sorted chronologically, checked for duplicate
timestamps, and aligned to a regular five-minute calendar. The
sample begins at 00:00 UTC on January 1, 2022 and ends at
23:55 UTC on December 31, 2025.

\subsection{Five-Minute Returns}

Let \(P_{t,j}\) denote the closing price of the \(j\)-th
five-minute interval within hour \(t\). The corresponding log
return is

\begin{equation}
r_{t,j}
=
\log P_{t,j}
-
\log P_{t,j-1}.
\end{equation}

Each return is timestamped at the end of its five-minute
interval and assigned to the hour in which the price change is
completed. The sample contains 420,767 valid returns; the first
observation is unavailable because no preceding closing price
exists.

\subsection{Hourly Realized Variance}

Realized variance aggregates high-frequency squared returns to
measure price variation over a specified interval
\citep{andersen2003}. The same approach has been applied to
Bitcoin volatility forecasting using intraday data
\citep{bergsli2022}. Let \(\mathcal{J}_t\) denote the twelve
five-minute returns completed during hour \(t\). Hourly
realized variance is

\begin{equation}
RV_t
=
\sum_{j \in \mathcal{J}_t}
r_{t,j}^{2}.
\label{eq:realized-variance}
\end{equation}

This measure incorporates the intrahour price path and is more
informative about short-term variation than a single hourly
return.

The aggregation produces 35,063 complete hourly bins. Because
realized variance is positive and strongly right-skewed, the
forecasting variable is defined as

\begin{equation}
y_t
=
\log(RV_t).
\label{eq:log-rv}
\end{equation}

No additive constant is introduced before taking logarithms.
Hours without a valid volatility measure are retained on the
hourly calendar and marked as missing.

\subsection{Data Quality Control}

A zero-activity bar is defined as a five-minute interval with
identical opening, high, low, and closing prices, zero trading
volume, and zero recorded trades. The data contain 55 such
bars, grouped into eight runs and affecting 11 hourly
observations.

Affected hours remain on the original calendar, but their
forecasting targets are set to missing. Lags and rolling
variables are constructed before any model-ready filtering.
This procedure preserves the temporal meaning of each lag and
prevents observations separated by missing hours from being
treated as adjacent.

After quality control, 35,052 hourly observations have valid
realized variance and log realized variance.

\subsection{Multi-Horizon Volatility Features}

The HAR framework represents volatility persistence through
aggregated components measured over different horizons
\citep{corsi2009}. This study adapts the framework to hourly
forecasting using one-hour, six-hour, daily, and weekly
components:

\[
\begin{aligned}
HAR^{(1)}_t
&=
y_t,\\
HAR^{(6)}_t
&=
\frac{1}{6}
\sum_{i=0}^{5}
y_{t-i},\\
HAR^{(24)}_t
&=
\frac{1}{24}
\sum_{i=0}^{23}
y_{t-i},\\
HAR^{(168)}_t
&=
\frac{1}{168}
\sum_{i=0}^{167}
y_{t-i}.
\end{aligned}
\]

These variables summarize current-hour, recent intraday,
daily, and weekly volatility conditions. All rolling averages
are calculated on the complete hourly calendar. A feature is
marked as missing whenever its required window contains an
invalid observation.

\subsection{Downside Realized Semivariance}

Realized semivariance separates total realized variation into
components generated by negative and positive intraday returns
\citep{barndorffnielsen2010}. Hourly downside realized
semivariance is defined as

\begin{equation}
RS^{-}_t
=
\sum_{j \in \mathcal{J}_t}
r_{t,j}^{2}
\mathbb{I}(r_{t,j}<0),
\label{eq:downside-semivariance}
\end{equation}

where \(\mathbb{I}(\cdot)\) is the indicator function.

Because the raw measure is highly right-skewed, it is
transformed using a training-based scale:

\begin{equation}
DS_t
=
\log
\left(
1+
\frac{RS^{-}_t}{s_{\mathrm{train}}}
\right),
\label{eq:downside-transform}
\end{equation}

where \(s_{\mathrm{train}}\) is the median of positive
downside-semivariance observations in the model-ready training
sample. Its estimated value is
\(s_{\mathrm{train}}=5.2624359\times10^{-6}\). The same value
is applied to the validation and test samples.

\subsection{Trading Activity}

Hourly trading activity is measured by the total number of
trades across the twelve five-minute intervals within each
hour. Because the distribution is strongly right-skewed, the
transformed measure is
\(Trades_t=\log(1+N_t)\), where \(N_t\) is the hourly trade
count.

Trading activity provides a proxy for information arrival and
market participation. The mixture-of-distributions framework
links trading intensity and price variability through common
information flows \citep{clark1973}, while empirical evidence
documents a close association between transaction counts and
volatility \citep{jones1994}. The variable is included to test
whether current market activity predicts next-hour volatility
after controlling for recent realized variance.

\subsection{Aggressive-Buying Proxy}

The Binance data report taker-buy base volume, representing
volume initiated by market buy orders. The aggressive-buying
proxy is

\begin{equation}
Buy_t
=
2
\left(
\frac{TBV_t}{V_t}
\right)
-1,
\label{eq:buy-pressure}
\end{equation}

where \(TBV_t\) is hourly taker-buy base volume and \(V_t\) is
total hourly base volume.

The variable is approximately bounded between \(-1\) and
\(1\). Positive values indicate that aggressive buying
accounts for more than half of total volume, while negative
values indicate stronger aggressive selling. Its sample mean
is close to zero.

\subsection{Feature Standardization}

The HAR-X predictor set comprises

\[
\left[
HAR^{(1)}_t,\,
HAR^{(6)}_t,\,
HAR^{(24)}_t,\,
HAR^{(168)}_t,\,
DS_t,\,
Trades_t,\,
Buy_t
\right].
\]

Each predictor is standardized using the mean and standard
deviation estimated from the model-ready training sample:

\begin{equation}
z_{k,t}
=
\frac{x_{k,t}-\bar{x}_{k,\mathrm{train}}}
{s_{k,\mathrm{train}}},
\label{eq:standardization}
\end{equation}

where \(x_{k,t}\) is predictor \(k\),
\(\bar{x}_{k,\mathrm{train}}\) is its training-sample mean,
and \(s_{k,\mathrm{train}}\) is its training-sample standard
deviation.

The transformation parameters remain fixed during validation
and test evaluation, preventing information from either period
from entering the preprocessing stage.

\subsection{Final Modeling Sample}

The complete hourly calendar contains 35,063 rows. Requiring a
valid one-hour-ahead target and complete values for all HAR-X
features yields 33,705 observations for the common model
comparison.

\begin{table}[H]
\centering
\caption{Final model-ready sample by period}
\label{tab:model-sample}
\begin{tabular}{lrr}
\toprule
Period & Calendar observations & Model-ready observations \\
\midrule
Training: 2022--2023 & 17,518 & 16,673 \\
Validation: 2024 & 8,784 & 8,610 \\
Test: 2025 & 8,760 & 8,422 \\
\midrule
Total used for modeling & 35,062 & 33,705 \\
\bottomrule
\end{tabular}
\end{table}

The table excludes the final forecast origin because its
one-hour-ahead target extends beyond the sample. That timestamp
remains in the complete hourly calendar used for feature
construction, which accounts for the difference between
35,063 calendar rows and the table total of 35,062.